\section{Abstract}
Digital Engineering has transformed how the Department of War, NASA, and the aerospace industry sustain complex systems. Its four pillars are Model-Based Systems Engineering, digital threads, digital twins, and Product Lifecycle Management, and together they are credited with gains in mission assurance and lifecycle governance. These same methods remain virtually untested in enterprise Information Technology and Information Assurance, even though compliance regimes such as the NIST Risk Management Framework and the DoW Cybersecurity Maturity Model Certification strain the documentation and traceability practices those fields depend on. A structured literature review found almost no research applying Digital Engineering to either domain, which leaves the question of whether these methods could transfer resting on argument rather than on evidence.

This study therefore asked the professionals themselves. A quantitative, cross-sectional survey measured what practitioners in Information Technology and Information Assurance know of these capabilities and what they believe those capabilities are worth. It drew 1{,}411 responses, roughly three times the planned target, from professional networks weighted toward defense and aerospace and toward earlier-career respondents.

Familiarity with Digital Engineering as a named discipline proved low. Seventy-eight percent of respondents placed themselves in the two lowest familiarity categories, and only 22\% reported any professional exposure within the previous two years. Once the survey described the underlying capabilities in plain functional language, however, agreement became near-universal at 94\%, and perceived value stayed high across the applicability, information technology, and security domains, ranging from 82\% to 88\%. Two subgroup patterns held up under scrutiny. Information Technology respondents outscored Information Assurance respondents by a wide margin, with effect sizes running from 1.64 to 1.94, and respondents with eleven to fifteen years of experience were more reserved than either their junior or their senior colleagues, a difference that survived adjustment for professional field and familiarity.

Two boundaries qualify these results. Recruitment leaned toward the sectors where Digital Engineering is already practiced, so the estimates should be read as upper-leaning rather than as measurements of the profession at large. The study also measured only two of the barriers that could stand in the way of adoption. Between those two, awareness and exposure bind more tightly than doubt about value, although a cross-sectional design cannot confirm that direction, and cost, tooling maturity, and organizational readiness went unmeasured.

The work contributes a first empirical baseline for a population that had not previously been measured, together with an open-source screening method, validated against synthetic data, for detecting coordinated or automated responses in anonymous surveys. Its wider significance turns on who carries the burden of compliance. That burden falls hardest on organizations with the least capacity to absorb it, and the professional receptiveness measured here is a precondition for the lower-cost implementations those organizations would need.
\clearpage
